\section{Comparison maps into a sphere}
\label{sec:maps}

We seek a sphere-valued map with a controlled differential and a
compatible Hessian bound for its weight. These two estimates will combine
in the hypersurface Sobolev energy. A second step chooses the map's
parameters to balance a given measure.

Fix $b\in\{0,1\}$ and a complete simply connected manifold $(M^n,g)$
with $\Sec_g\le-b^2$. We use the model functions $\sn_b$ and $\cs_b$
from Section~\ref{sec:introduction}; thus
\[
 \sn_b'=\cs_b,\qquad \cs_b'=b^2\sn_b,
 \qquad \cs_b^2-b^2\sn_b^2=1.
\]
The Cartan--Hadamard theorem makes $\exp_p:T_pM\to M$ a diffeomorphism
for every $p$; see \cite[Theorem~6.2.2]{PetersenManuscript}.
For $x\ne p$, put $r=d(p,x)$ and $\theta=\exp_p^{-1}(x)/r$.
The sphere $\mathbb S(T_pM\oplus\R)$ has its round metric induced by
$g_p$ and the usual metric on $\R$.
Write $g_{\mathbb S_p^{n-1}}$ for the round metric on the unit sphere
in $T_pM$.

In the radial coordinate $\sn_b(r/2)/\cs_b(r/2)$, the model polar
metric $dr^2+\sn_b(r)^2g_{\mathbb S_p^{n-1}}$ is conformal to the
Euclidean polar metric. We dilate this coordinate by $a>0$ and apply
inverse stereographic projection, obtaining the map below and its model
conformal factor $f_{p,a}$:
\begin{equation}\label{eq:map-definition}
 \begin{gathered}
 z=a\frac{\sn_b(r/2)}{\cs_b(r/2)},\qquad
 \Phi_{p,a}(x)=\left(\frac{2z}{1+z^2}\theta,
                          \frac{1-z^2}{1+z^2}\right),\\
 f_{p,a}(x)=e^{u_{p,a}(x)}
 =\frac{a}{\cs_b(r/2)^2+a^2\sn_b(r/2)^2}.
 \end{gathered}
\end{equation}
At the pole set $\Phi_{p,a}(p)=(0,1)$ and $f_{p,a}(p)=a$.
In particular the zero-curvature formulas are explicitly
$z=ar/2$ and $f_{p,a}=a/(1+a^2r^2/4)$.

\begin{lemma}[Comparison maps]\label{lem:comparison-maps}
The map $\Phi_{p,a}$ and the positive function $f_{p,a}$ are smooth on
$M$. For every $X\in T_xM$,
\begin{equation}\label{eq:map-comparison}
 \begin{aligned}
 |d\Phi_{p,a}(X)|&\le f_{p,a}|X|,\\
 \Hess u_{p,a}&\le du_{p,a}\otimes du_{p,a}
 -\frac12\bigl(b^2+f_{p,a}^2+|\nabla u_{p,a}|^2\bigr)g.
 \end{aligned}
\end{equation}
The second inequality is one of quadratic forms. The functions depend
smoothly on $(p,a,x)$. For fixed $o\in M$, parallel transport of the
first component of $\Phi_{p,a}$ from $p$ to $o$ along the unique geodesic
also depends smoothly on $(p,a,x)$.
\end{lemma}

\begin{proof}
The two components of $\Phi_{p,a}$ have squared norms summing to one.
In normal coordinates $y=\exp_p^{-1}(x)$, its first component is
$2z\,y/(r(1+z^2))$. The coefficient, the last component, and $f_{p,a}$
are smooth functions of $r^2$ at zero. Their denominators are positive.
This proves smoothness at the pole, including that of $u=\log f$.
The map $(p,y)\mapsto(p,\exp_p y)$ from $TM$ to $M\times M$ is a
bijective local diffeomorphism, hence a diffeomorphism. Its inverse
and the same even radial expressions prove joint smoothness.
The paths $t\mapsto\exp_o(t\exp_o^{-1}p)$ vary smoothly with $p$;
smooth dependence for the parallel-transport equation proves the
last assertion.

We next establish the one-sided distance comparison needed below:
\begin{equation}\label{eq:distance-hessian-comparison}
 \mathcal H:=\Hess r-\frac{\cs_b(r)}{\sn_b(r)}
                         (g-dr\otimes dr)\ge0\qquad(r>0).
\end{equation}
Indeed, for a vector $X$ perpendicular to a unit radial geodesic of
length $r$, the Jacobi field with endpoint values $0,X$ gives
\[
 \Hess r(X,X)
 =\int_0^r\bigl(|D_tJ|^2-\langle R(J,\dot\gamma)\dot\gamma,J\rangle\bigr)\,dt
 \ge\int_0^r(|D_tJ|^2+b^2|J|^2)\,dt.
\]
In a parallel trivialization the last energy is minimized by
$J_0(t)=\sn_b(t)X/\sn_b(r)$ and equals
$\cs_b(r)|X|^2/\sn_b(r)$. To verify the minimum, write $J=J_0+V$;
integration by parts makes the cross term zero because $V$ vanishes
at the endpoints and $J_0''=b^2J_0$. The remaining energy of $V$ is
nonnegative. The radial and mixed components of $\Hess r$ vanish,
proving \eqref{eq:distance-hessian-comparison} without a lower
curvature bound.

For an angular variation field $J$ along a radial geodesic, Gauss's
lemma and commutation of the variation fields give
\[
 \frac d{dr}|J|^2=2\Hess r(J,J)
 \ge2\frac{\cs_b(r)}{\sn_b(r)}|J|^2.
\]
Thus $|J|^2/\sn_b(r)^2$ is nondecreasing, with limit equal to the
initial angular squared norm at zero. Applying this to every angular
vector proves the polar metric comparison
\begin{equation}\label{eq:polar-comparison}
 g\ge dr^2+\sn_b(r)^2g_{\mathbb S_p^{n-1}}.
\end{equation}
The round stereographic metric is
$4(dz^2+z^2g_{\mathbb S_p^{n-1}})/(1+z^2)^2$.
The identities
\[
 \frac{2z'}{1+z^2}=f,\qquad
 \frac{2z}{1+z^2}=f\sn_b(r)
\]
therefore give
\begin{equation}\label{eq:map-pullback}
 \Phi_{p,a}^*g_{\mathbb S(T_pM\oplus\R)}
 =f^2\bigl(dr^2+\sn_b(r)^2g_{\mathbb S_p^{n-1}}\bigr)\le f^2g.
\end{equation}
This proves the differential estimate off the pole.

For the Hessian calculation write
\[
 A=a^2+b^2,\qquad
 D=\cs_b(r/2)^2+a^2\sn_b(r/2)^2=1+A\sn_b(r/2)^2.
\]
Then $f=a/D$, $D'=A\sn_b(r)/2$, and $D''=A\cs_b(r)/2$.
The double-angle identities give
\[
 b^2D^2+a^2+\frac{A^2}{4}\sn_b(r)^2=A\cs_b(r)D.
\]
For example, substituting $t=\sn_b(r/2)^2$ makes both sides
$A+(A^2+2b^2A)t+2b^2A^2t^2$. Consequently
\begin{equation}\label{eq:radial-identities}
 \begin{gathered}
 u'=-\frac{A\sn_b(r)}{2D}\le0,\qquad
 u''=(u')^2+u'\frac{\cs_b(r)}{\sn_b(r)},\\
 u'\frac{\cs_b(r)}{\sn_b(r)}
 =-\frac12\bigl(b^2+f^2+(u')^2\bigr).
 \end{gathered}
\end{equation}
Combining these identities with \eqref{eq:distance-hessian-comparison},
\[
 \Hess u
 =du\otimes du-\frac12\bigl(b^2+f^2+|\nabla u|^2\bigr)g
       +u'\mathcal H.
\]
The last term is negative semidefinite. Finally, the pole expansions
\[
 d\Phi_{p,a}|_p(X)=(aX,0),\qquad
 u=\log a-\frac{a^2+b^2}{4}r^2+O(r^4)
\]
show that both estimates in \eqref{eq:map-comparison} also hold at $p$.
\end{proof}

The zero-curvature case also gives three comparisons that will be used
in the equality argument. They express, respectively, convexity of
squared distance, expansion by the exponential map, and contraction
by its inverse.

\begin{corollary}[Distance and metric comparison]\label{lem:distance-comparison}
Let $(M^n,g)$ be complete and simply connected with $\Sec_g\le0$.
For $p\in M$, write $\log_p=\exp_p^{-1}$, $r_p(x)=d(p,x)$, and
$D_p=r_p^2/2$. Then
\begin{equation}\label{eq:distance-metric-comparison}
 \Hess_gD_p\ge g,\qquad (\exp_p)^*g\ge g_p,
 \qquad |d\log_p(X)|_{g_p}\le |X|_g.
\end{equation}
Here $g_p$ on $T_pM$ denotes its constant Euclidean metric.
\end{corollary}

\begin{proof}
With $b=0$, \eqref{eq:distance-hessian-comparison} gives
\[
 \Hess D_p=dr_p\otimes dr_p+r_p\Hess r_p\ge g
\]
away from $p$; smoothness of squared distance and normal coordinates
give equality at $p$. The $b=0$ case of
\eqref{eq:polar-comparison} is exactly $(\exp_p)^*g\ge g_p$,
including at the origin by continuity. Applying this differential
inequality to the inverse map gives the last assertion.
\end{proof}

The map parameters have supplied the geometric estimates. We next
choose the center and scale so that the coordinate functions have zero
mean for a prescribed measure. This use of spherical centering is
related to the conformal variational argument of Li and Yau
\cite[pp.~274--275]{LiYau1982}; the required centering statement is
proved below for these comparison maps.

\begin{lemma}[Balancing a nonatomic measure]\label{lem:balancing}
Let $\mu$ be a compactly supported Borel probability measure on $M$
with $\mu(\{x\})=0$ for every $x$. If $\supp\mu\subset B_D(o)$,
there are $p\in B_D(o)$ and finite $a>0$ such that
\begin{equation}\label{eq:balanced-map}
 \int_M\Phi_{p,a}(x)\,d\mu(x)=0\quad\text{in }T_pM\oplus\R.
\end{equation}
\end{lemma}

\begin{proof}
Write $K=\supp\mu$, parameterize $p$ by $\xi=\exp_o^{-1}p$, and let
$P_{p\to o}$ be radial parallel transport. Define
\[
 F(\xi,a)=(F_1,F_2)
 =\int_M(P_{p\to o}\oplus\mathrm{id}_{\R})\Phi_{p,a}(x)\,d\mu(x).
\]
It is continuous on every cylinder
$\mathcal C=\{\xi\in T_oM:|\xi|\le D\}\times[a_0,a_1]$
with $0<a_0<a_1<\infty$: the integrand is jointly continuous and has
norm one. We choose the two horizontal faces so that $F$ points
strictly inward on the entire boundary.

First suppose $|\xi|=D$. Put $\rho=d(o,\cdot)$. Radial parallel
transport sends $\xi$ to $D\nabla\rho(p)$. Convexity of $\rho$ along
the geodesic from $p$ to $x\in K$ gives
\[
 \langle\xi,P_{p\to o}\theta_p(x)\rangle
 \le\frac{D(\rho(x)-D)}{d(p,x)}<0.
\]
Here $p\notin K$; the derivative is taken only at $p\ne o$, so a
geodesic passing through $o$ causes no difficulty. The first component
of $\Phi_{p,a}$ is a positive multiple of $\theta_p(x)$ on this face.
Hence
\begin{equation}\label{eq:balance-side}
 \langle\xi,F_1(\xi,a)\rangle<0
 \qquad(|\xi|=D,\ a>0).
\end{equation}

Set $t_b(r)=\sn_b(r/2)/\cs_b(r/2)$, an increasing positive function
for $r>0$. Choose $a_0>0$ so that $a_0t_b(2D)<1$.
Since $d(p,x)<2D$ for $|\xi|\le D$ and $x\in K$, we obtain
\begin{equation}\label{eq:balance-bottom}
 F_2(\xi,a_0)\ge
 \frac{1-a_0^2t_b(2D)^2}{1+a_0^2t_b(2D)^2}>0.
\end{equation}

We also have
\begin{equation}\label{eq:uniform-atomless}
 \lim_{\delta\downarrow0}\sup_{p\in\overline B_D(o)}
                            \mu(B_\delta(p))=0.
\end{equation}
Otherwise there would be $\varepsilon>0$, $\delta_j\to0$, and centers
$p_j$ with $\mu(B_{\delta_j}(p_j))\ge\varepsilon$. The closed ball is
compact by completeness, so a subsequence of centers converges to
$p_\infty$. Every $B_s(p_\infty)$ then has mass at least $\varepsilon$.
Continuity of a finite measure from above gives an atom at $p_\infty$,
a contradiction.

Choose $\delta>0$ so that the supremum in
\eqref{eq:uniform-atomless} is less than $1/4$, and choose finite
$a_1>a_0$ with $a_1t_b(\delta)\ge\sqrt3$.
The last coordinate of $\Phi_{p,a_1}$ is at most $-1/2$ outside
$B_\delta(p)$ and at most one everywhere. Thus
\begin{equation}\label{eq:balance-top}
 F_2(\xi,a_1)\le-\frac18<0\qquad(|\xi|\le D).
\end{equation}

Let $\Pi_{\mathcal C}$ be Euclidean nearest-point projection onto the
compact convex cylinder. Brouwer's theorem gives a fixed point of
$T(\zeta)=\Pi_{\mathcal C}(\zeta+F(\zeta))$. The projection inequality
at that point is
\begin{equation}\label{eq:projection-fixed-point}
 \langle F(\zeta),y-\zeta\rangle\le0
                         \qquad(y\in\mathcal C).
\end{equation}
On the side face, choosing $y=((1-t)\xi,a)$ with small $t>0$ contradicts
\eqref{eq:balance-side}. On the lower or upper face, moving the last
coordinate inward contradicts \eqref{eq:balance-bottom} or
\eqref{eq:balance-top}. These choices remain admissible at corners.
The fixed point is therefore interior. Taking both signs of every
small displacement in \eqref{eq:projection-fixed-point} gives $F=0$.
Parallel transport is an isomorphism, which proves
\eqref{eq:balanced-map} with the asserted parameters.
\end{proof}
